\documentclass[12pt,titlepage]{article}

\usepackage{amsmath}
\usepackage{xcolor}
\usepackage{mathtools}
\usepackage{indentfirst}
\usepackage{hhline}

\title{\bf Largest $m$-cube in an $n$-cube: \\
       Families
      }

\author{Greg Huber \\
        Terry J. Ligocki
       }

\begin{document}

\maketitle

\begin{abstract}
In looking at the largest $m$-cube in an $n$-cube problem, $(m,n)$ problem,
various patterns occur in subsets of $(m,n)$ where m is a function of $n$,
$(g(n),n)$, or $n$ is a function of $m$, $(m,g(m))$.  We call these "families".
In this document we list all the patterns/families we have found along with
any formulas relating members of a family.  We look at both the maximal ratio
of the side length of the $m$-cube to the side length of the $n$-cube and the
orthogonal matrix that defines the imbedding of the $m$-cube within the
$n$-cube.
\end{abstract}

\section{Patterns and Families}

\subsection{$(1,n)$ Family (Provable)}
\label{sec:1_n}
The ratio of the side lengths is $\sqrt{n}$ which is achieved when the 1-cube
(a line segement) is aligned any two corners that are reflections through the
origin.  The distance between these corners is the largest diameter of the
cube so this is largest value the $1$-cube could have within the $n$-cube.

The orthogonal matrix has $1/\sqrt{n}$ for each entry in the first column.
The rest of the columns are unrestricted except that the entire matrix must
be orthogonal.
\begin{equation*}
l(1,n) = \sqrt{n}
\end{equation*}
\begin{equation*}
M(1,n) =
\begin{bmatrix}
1/\sqrt{n} & \ldots \\
\vdots     & \vdots \\
1/\sqrt{n} & \ldots
\end{bmatrix}
\end{equation*}

\subsection{$(2,2n)$ Family (Provable)}
\label{sec:2_2n}
This reduces to the $(1,n)$ case via the general conjecture $l(am,an) = l(m,n)$.
This conjecture also leads to orthogonal matrix where the first m columns of
the $(m,n)$ orthogonal matrix are copied into the $(am,an)$ orthogonal matrix
starting at index $[im,in]$ for integer $i$ in $\{0, \ldots, a-1\}$.  All the
other entries in the first $am$ columns are zeros.
\begin{equation*}
l(2,2n) = \sqrt{n}
\end{equation*}
\begin{equation*}
M(2,2n) =
\begin{bmatrix}
1/\sqrt{n} & 0          & \ldots \\
\vdots     & \vdots     & \vdots \\
1/\sqrt{n} & 0          & \ldots \\
0          & 1/\sqrt{n} & \ldots \\
\vdots     & \vdots     & \vdots \\
0          & 1/\sqrt{n} & \ldots \\
\end{bmatrix}
\end{equation*}

\subsection{$(2,2n+1)$ Family (Provable)}
\label{sec:2_2n+1}
\textcolor{red}{
Add text about this case.
}
\begin{equation*}
l(2,2n+1) = \sqrt{n + 1/8}
\end{equation*}
\begin{equation*}
M(2,2n+1) =
\begin{bmatrix}
\textcolor{red}{Update}
\end{bmatrix}
\end{equation*}

\subsection{$(am,an)$ is equivalent to $(m,n)$ (Conjectured)}
\textcolor{red}{
Add text about this case.
}
\begin{equation*}
l(am,an) = l(m,n)
\end{equation*}
\begin{equation*}
M(am,an) =
\begin{bmatrix}
M_m(m,n) & 0      & \ldots & 0        & \ldots \\
0        & \ddots &        & \vdots   & \vdots \\
\vdots   &        & \ddots & 0        & \vdots \\
0        & \ldots & 0      & M_m(m,n) & \ldots
\end{bmatrix}
\end{equation*}
Where $M_m(m,n)$ is the first $m$ columns of the matrix $M(m,n)$ for $(m,n)$.

\subsection{$(n-1,n)$ Family (Computed)}
\textcolor{red}{
Describe technique for computing these.  Include a table of computed
$l(n-1,n) - 1$ and the ratio below.
}
\begin{equation*}
l(n-1,n) = \textcolor{red}{Update}
\end{equation*}
\begin{equation*}
M(n-1,n) =
\begin{bmatrix}
\textcolor{red}{Update}
\end{bmatrix}
\end{equation*}
\begin{equation*}
\lim_{n \to \infty}\frac{l(n-1,n) - 1}{l(n,n+1) - 1} = 9
~(\textcolor{red}{conjectured})
\end{equation*}

\subsection{$(n-2,n)$ Family (Computed)}
\textcolor{red}{
Describe technique for computing these.  Include a table of computed
$l(n-2,n)$.
}
\begin{equation*}
l(n-1,n) = \textcolor{red}{Update}
\end{equation*}
\begin{equation*}
M(n-1,n) =
\begin{bmatrix}
\textcolor{red}{Update}
\end{bmatrix}
\end{equation*}

\begin{table}
\begin{center}
\begin{tabular}{| c | c | c | c |}
\hline
$n-1$ & $n$ & $l(n-1,n)-1$    & $\frac{l(n-1,n)-1}{l(n,n+1)-1}$ \\
\hhline{|=|=|=|=|}
  1 &   2 & 4.1421356237e-01 & N/A \\
\hline
  2 &   3 & 6.0660171780e-02 & 6.8284271247e+00 \\
\hline
  3 &   4 & 7.4347568843e-03 & 8.1589987035e+00 \\
\hline
  4 &   5 & 8.3944685935e-04 & 8.8567332184e+00 \\
\hline
  5 &   6 & 9.3473576520e-05 & 8.9805792247e+00 \\
\hline
  6 &   7 & 1.0388837407e-05 & 8.9975011507e+00 \\
\hline
  7 &   8 & 1.1543556344e-06 & 8.9996852762e+00 \\
\hline
  8 &   9 & 1.2826229404e-07 & 8.9999609242e+00 \\
\hline
  9 &  10 & 1.4251373602e-08 & 8.9999952022e+00 \\
\hline
 10 &  11 & 1.5834860585e-09 & 8.9999994162e+00 \\
\hline
 11 &  12 & 1.7594289676e-10 & 8.9999999295e+00 \\
\hline
 12 &  13 & 1.9549210770e-11 & 8.9999999915e+00 \\
\hline
 13 &  14 & 2.1721345302e-12 & 8.9999999990e+00 \\
\hline
 14 &  15 & 2.4134828114e-13 & 8.9999999999e+00 \\
\hline
 15 &  16 & 2.6816475682e-14 & 9.0000000000e+00 \\
\hline
\end{tabular}
\end{center}
\vspace{-0.2in}
\caption{\textcolor{red}{This explains the table.}}
\end{table}

\subsection{$(m,2m-1)$ Family (Computed/Conjectured)}
\label{sec:m_2m-1}
\textcolor{red}{
Add text about this case.
}
\begin{equation*}
l(m,2m-1) = \left(\frac{1 + \sqrt{2m}}{2 + \sqrt{2m}}\right) \sqrt{2}
\end{equation*}
\begin{equation*}
M(m,2m-1) =
\begin{bmatrix*}[r]
\alpha & 0      & \ldots & 0      & -\beta  & \ldots \\
\alpha & 0      & \ldots & 0      &  \beta  & \ldots \\
0      & \ddots &        & \vdots & \vdots   & \vdots \\
\vdots &        & \ddots & 0      & \vdots   & \vdots \\
0      & \ldots & 0      & \alpha & -\beta  & \ldots \\
0      & \ldots & 0      & \alpha &  \beta  & \ldots \\
0      & \ldots & \ldots & 0      &  \gamma & \ldots
\end{bmatrix*}
\end{equation*}
Where:
\begin{align*}
\alpha &= \frac{\sqrt{2}}{2}                          \\
\beta  &= \left(\frac{1}{1 + \sqrt{2m}}\right) \alpha \\
\gamma &= \left(2 + \sqrt{2m}\right) \beta 
        = \left(\frac{2 + \sqrt{2m}}{1 + \sqrt{2m}}\right) \alpha
\end{align*}

\subsection{$m$ is fixed and $n \bmod m$ is a contant}

In this case, it appears that when $n$ is ``large enough'', $l(m,n)$ and
$M(m,n)$ become a family with a regular pattern of extension.  Here are some
examples:

\subsubsection{(m,nm) --- $n \bmod m = mn \bmod m = 0$}
This reduces to the (1,n) cases, see section \ref{sec:1_n}.

\subsubsection{(2,2n+1) --- $n \bmod m = 2n+1 \bmod 2 = 1$}
Here there is a recursive definition of $l(2,2n+1)$ and $M(2,2n+1)$:
\begin{align*}
l(2,3)        &= \sqrt{9/8} 
               = \frac{3}{4} \sqrt{2} 
               \approx 1.0606602        \\
l(2,2(n+1)+1) &= \sqrt{l(2,2n+1)^2 + 1}
\end{align*}
\begin{align*}
M(2,3) &=
\begin{bmatrix*}[r]
\sqrt{2}/2 & -\sqrt{2}/6 &  2/3 \\
\sqrt{2}/2 &  \sqrt{2}/6 & -2/3 \\
0          & 2\sqrt{2}/3 &  1/3
\end{bmatrix*} \\
M(2,2(n+1)+1) &=
\begin{bmatrix*}[c]
\rho~M_2(2,2n+1)  & \ldots \\
l(2,2(n+1)+1)~I_2 & \ldots
\end{bmatrix*}
\end{align*}
Where, $M_2(2,2n+1)$ is the first two columns of $M(2,2n+1)$,
$\rho = l(2,2(n+1)+1)/l(2,2n+1)$, and $I_2$ is the 2 by 2 identity matrix.

See section \ref{sec:2_2n+1} for an alternative presentation of this family.

\subsubsection{(3,3n+1) --- $n \bmod m = 3n+1 \bmod 3 = 1$}
There are recursive definitions of $l(3,3n+1)$ and $M(3,3n+1)$ once $m =
3$:
\begin{align*}
l(3,10)       &\approx 1.7416566        \\
l(3,3(n+1)+1) &= \sqrt{l(3,3n+1)^2 + 1}
\end{align*}
\begin{align*}
M(3,10) &=
\begin{bmatrix*}[c]
 -(1/l(3,10) - 1/2) & 0                 & 1/2 & \ldots \\
~~(1/l(3,10) - 1/2) & 0                 & 1/2 & \ldots \\
0                 &  -(1/l(3,10) - 1/2) & 1/2 & \ldots \\
0                 & ~~(1/l(3,10) - 1/2) & 1/2 & \ldots \\
1/l(3,10)           & 0                 & 0   & \ldots \\
0                 & 1/l(3,10)           & 0   & \ldots \\
1/l(3,10)           & 0                 & 0   & \ldots \\
0                 & 1/l(3,10)           & 0   & \ldots \\
1/l(3,10)           & 0                 & 0   & \ldots \\
0                 & 1/l(3,10)           & 0   & \ldots 
\end{bmatrix*} \\
M(3,3(n+1)+1) &=
\begin{bmatrix*}[c]
\rho~M_3(3,3n+1)  & \ldots \\
l(3,3(n+1)+1)~I_3 & \ldots
\end{bmatrix*}
\end{align*}
Where, $M_3(3,3n+1)$ is the first three columns of $M(3,3n+1)$,
$\rho = l(3,3(n+1)+1)/l(3,3n+1)$, and $I_3$ is the 3 by 3 identity matrix.

\subsubsection{(3,3n+2) --- $n \bmod m = 3n+2 \bmod 3 = 2$}
First, $l(3,5)$ and $M(3,5)$ are ``known'' because $(3,5) = (3,2*3-1)$ and
conjectured solutions for this family are given in section \ref{sec:m_2m-1}.
Here are recursive definitions of $l(3,3n+2)$ and $M(3,3n+2)$:
\begin{align*}
l(3,5)        &= \left(\frac{1+\sqrt{6}}{2+\sqrt{6}}\right) \sqrt{2}
               = \frac{(4-\sqrt{6}) \sqrt{2}}{2}
               \approx 1.0963763        \\
l(3,3(n+1)+2) &= \sqrt{l(3,3n+2)^2 + 1}
\end{align*}
\begin{align*}
M(3,5) &=
\begin{bmatrix*}[r]
\alpha & 0      & -\beta  & \ldots \\
\alpha & 0      &  \beta  & \ldots \\
0      & \alpha & -\beta  & \ldots \\
0      & \alpha &  \beta  & \ldots \\
0      & 0      &  \gamma & \ldots
\end{bmatrix*} \\
& \\
\alpha &= \sqrt{2}/2                \\
\beta  &= \sqrt{2}(\sqrt{6} - 1)/10 \\
\gamma &= \sqrt{2}(\sqrt{6} + 4)/10 \\
& \\
M(3,3(n+1)+2) &=
\begin{bmatrix*}[c]
\rho~M_3(3,3n+2)  & \ldots \\
l(3,3(n+1)+2)~I_3 & \ldots
\end{bmatrix*}
\end{align*}
Where, $M_3(3,3n+2)$ is the first three columns of $M(3,3n+2)$,
$\rho = l(3,3(n+1)+2)/l(3,3n+2)$, and $I_3$ is the 3 by 3 identity matrix.

\subsubsection{(4,4n+3) --- $n \bmod m = 4n+3 \bmod 4 = 3$}
First, $l(4,7)$ and $M(4,7)$ are ``known'' because $(4,7) = (4,2*4-1)$ and
conjectured solutions for this family are given in section \ref{sec:m_2m-1}.
Here are recursive definitions of $l(4,4n+3)$ and $M(4,4n+3)$:
\begin{align*}
l(4,7)        &= \left(\frac{1+\sqrt{8}}{2+\sqrt{8}}\right) \sqrt{2}
               = \frac{(3\sqrt{2}-2)}{2}
               \approx 1.1213203        \\
l(4,4(n+1)+3) &= \sqrt{l(4,4n+3)^2 + 1}
\end{align*}
\begin{align*}
M(4,7) &=
\begin{bmatrix*}[r]
\alpha & 0      & 0      & -\beta  & \ldots \\
\alpha & 0      & 0      &  \beta  & \ldots \\
0      & \alpha & 0      & -\beta  & \ldots \\
0      & \alpha & 0      &  \beta  & \ldots \\
0      & 0      & \alpha & -\beta  & \ldots \\
0      & 0      & \alpha &  \beta  & \ldots \\
0      & 0      & 0      &  \gamma & \ldots
\end{bmatrix*} \\
& \\
\alpha &= \sqrt{2}/2         \\
\beta  &= (4  - \sqrt{2})/14 \\
\gamma &= (3\sqrt{2} + 2)/7  \\
& \\
M(4,4(n+1)+3) &=
\begin{bmatrix*}[c]
\rho~M_4(4,4n+3)  & \ldots \\
l(4,4(n+1)+3)~I_4 & \ldots
\end{bmatrix*}
\end{align*}
Where, $M_4(4,4n+3)$ is the first four columns of $M(4,4n+3)$,
$\rho = l(4,4(n+1)+3)/l(4,4n+3)$, and $I_4$ is the 4 by 4 identity matrix.

\subsubsection{General extensions of $(m,n)$ to $(m,n+km)$}
In general, if one has a candidate solution for $(m,n)$, this can be extended
to $(m,n+km)$.  This is done recursively using that following constructions
starting with $i = 0$:
\begin{equation*}
l(m,n+(i+1)m) = \sqrt{l(m,n+im)^2 + 1}
\end{equation*}
\begin{equation*}
M(m,n+(i+1)m) =
\begin{bmatrix*}[c]
\rho~M_m(m,n+im)  & \ldots \\
l(m,n+(i+1)m)~I_m & \ldots
\end{bmatrix*}
\end{equation*}
Where, $M_m(m,n+im)$ is the first $m$ columns of $M(m,n+im)$,
$\rho = l(m,n+(i+1)m)/l(m,n+im)$, and $I_m$ is the $m$ by $m$ identity matrix.

These extensions aren't guaranteed to be optimal even if the candidate
solution is but they form a lower bound for $l(m,n+km)$.  In addition, they do
appear to be optimal in some of the cases we have computed.

\end{document}
